\setcounter{section}{1}
\section{Some theories of semirings with fixed points}
\label{sec:prelims}

Let $\Alphabet$ be a finite alphabet and $\Var$ be a countably infinite set of variables disjoint from $\Alphabet$. 
We consider \defname{$\mu$-expressions} $e,f,$ etc.\ generated by the following grammar:
\[
e,f,\dots \quad ::= \quad
0 \quad \mid \quad 1 \quad \mid \quad a \in \Alphabet \quad \mid \quad x\in\Var \quad \mid \quad e+f \quad \mid \quad e\cdot f \quad \mid \quad \mu x e
\]
We usually just write $ef$ instead of $e\cdot f$.

Structures interpreting $\mu$-expressions, formally, must be defined appropriately taking account of variable binding and compatible with substitution. 
A formal exposition can be found in e.g.\ \cite{EsikLeiss05,Leiss16}, so here we give only examples:

\begin{example}
    [Languages]
    \label{language-model}
    Temporarily expand the syntax of expressions to include each language $A\subseteq \Alphabet^*$ as a new constant symbol.
    We associate to each closed $\mu$-expression $e$ (of this expanded syntax) a language $\lang e \subseteq \Alphabet^*$ as follows:
    \[
    \begin{array}{r@{\ \df \ }l}
         \lang 0 & 0 \\
         \lang 1 & \{\epsilon\} \\
         \lang a & \{a\} \\
         \lang A & A\\
         \lang {e+f} & \lang e \cup \lang f \\
          \lang {ef} & \{ uv : u \in \lang e \ \& \ v \in \lang f\} \\
         \noalign{\medskip}
          \lang {\mu x e(x)} & \bigcap\{ A : \lang {e(A)} \subseteq A \}
    \end{array}
    \]
    % \[
    % \lang {\mu x e(x)} \quad \df \quad \bigcap\{ A : \lang {e(A)} \subseteq A \}
    % \]
    % \begin{itemize}
    %     \item $\lang 0 \df 0$
    %     \item $\lang 1 \df \{\epsilon\}$
    %     \item $\lang a \df \{a\}$
    %     \item $\lang A \df A$
    %     \item $\lang{e+f} \df \lang e \cup \lang f$
    %     \item $\lang {ef} \df \{uv : u \in \lang e \ \& \ v \in \lang f\}$
    %     \item $\lang {\mu x e(x)} \df \bigcap\{ A : \lang {e(A)} \subseteq A \}$
    % \end{itemize}
\end{example}
Every symbol is interpreted by $\lang\cdot$ as a monotone operation, with respect to $(\pow {\Alphabet^*}, \subseteq)$.
The \emph{Knaster-Tarski} fixed point theorem tells us that monotone operations on complete lattices have least fixed points,\footnote{Indeed, they further have a complete sublattice of fixed points, in particular including a greatest fixed point.} in particular with the least fixed point of $A \mapsto e(A)$ given by $\lang e $ defined above: the intersection of all \emph{prefixed points} of $A \mapsto e(A)$.

It is not hard to see that the image of $\lang \cdot$ on $\mu$-expressions is just the context-free languages (CFLs) \cite{Gruska71:cfls-by-mu-exprs}.
This structure will form a model of all theories considered herein, and will guide our main undecidability results.
Let us point out already the following consequence:

\begin{proposition}
[See, e.g., \cite{HopcroftUllman79:book,Sipser96:book}]
    Deciding $\lang e \subseteq_? \lang f$ is $\Pi^0_1$-complete.
\end{proposition}

\subsection{Semirings with fixed points}
For generality of our results, we shall work with a minimal theory of semirings with fixed points, before considering specific examples such as (idempotent) Park $\mu$-semirings and Conway $\mu$-semirings later.

\begin{definition}
[Basic axioms]
    The equational theory $\basicax$ of $\mu$-expressions is given by all the following axioms:
    \begin{enumerate}[(i)]
    \item $(0,1,+,\cdot) $ form a semiring:
    \[
    \begin{array}{r@{\ = \ }l}
         e + (f+g) & (e+f)+g \\
         e+f & f + e\\
         e+0 & e 
    \end{array}
    \quad
    \begin{array}{r@{\ = \ }l}
         e (fg) & (ef)g \\
         e1 & e \\
         1 e & e \\
         e0 & 0 \\
         0e & 0
    \end{array}
    \quad
    \begin{array}{r@{\ = \ }l}
         e(f+g) & ef + eg \\
         (e+f)g & eg + fg
    \end{array}
    \]
    % \item $\basicax$ is closed under \defname{$\mu$-congruence}: $\dfrac{e(x)=f(x)}{\mu x e(x)=\mu x f(x)}$.\footnote{This is a standard condition for equational theories with binders, in order to appropriately admit equational reasoning `within' terms. In particular, being an equational theory, equality in $\basicax$ satisfies the laws of a congruence.}
    % \item $e\leq f \iff e+f = f$;
    \item\label{item:fixed-point-axioms} $\mu x e(x)$ is a fixed point of $ x \mapsto e(x)$:
    \[
    \mu x e(x) = e (\mu x e(x))
    \]
    \end{enumerate}
\end{definition}

Note that $\basicax$ does \emph{not} insist that $\mu x e(x)$ is a \emph{least} fixed point, which would require non-equational axioms. 
Instead, notice that the axiomatisation above is finite and equational.


\begin{remark}
    [On $\mu$-congruence]
    \label{rem:mu-congruence}
    Note that $\basicax$ does not admit the \textbf{$\mu$-congruence} rule: $\dfrac{e(x)=f(x)}{\mu x e(x)=\mu x f(x)}$.
    This is often built into the definition of various classes of `$\mu$-semirings' in the literature \cite{EsikLeiss05,GHK13,Leiss16}, but we shall not exploit this here.
    This means that, even if two operators are (provably) equivalent, the fixed points named by their closure under $\mu$ may be different. 
    
    For instance we do not have in $\basicax$ that $\mu x (a+x) = \mu x (x+a)$.
    This is consistent with the fact that we are not axiomatising $\mu x e(x)$ to be a \emph{least} fixed point, but just \emph{some} fixed point.
    So, e.g., we can tweak the model $\Lang$ to satisfy $\mu x (a+x) =a $ but $\mu x (x+a) = a+b$.
    Of course only the former is the \emph{least} fixed point, but the resulting structure remains a model of $\basicax$.

    Construing expressions in a higher-order syntax, where $\mu$ is a higher-order constant rather than a binder, the lack of $\mu$-congruence may be seen as saying that $\mu$ does not behave \emph{extensionally}.
    In practice, in this work, this means we must be careful to not carry out reasoning \emph{under} a $\mu$ unless justified.
\end{remark}

\begin{remark}
    [Comparison to $\mu$-semirings]
    A \defname{$\mu$-semiring}, introduced in \cite{EsikLeiss05}, is defined just like $\basicax$ except:
    \begin{itemize}
        \item $\mu$-semirings must satisfy \emph{strong} $\mu$-congruence:
        $\forall x\,  e(x) = f(x) \limp \mu x e(x) = \mu x f(x)$; and,
        \item $\mu$-semirings do not require $\mu x e(x) $ to be a fixed point, i.e.\ axiom \eqref{item:fixed-point-axioms} is omitted.
    \end{itemize}
    Thus our system $\basicax$ is, strictly speaking, incomparable with general $\mu$-semirings, though it will still underlie extensions considered later.
     $\basicax$ should rather be considered a basic theory of `semirings with fixed points', namely the theory of those $\mu$-semirings in which $\mu x e(x)$ always denotes a fixed point of $x \mapsto e(x)$, albeit not necessarily the \emph{least} fixed point, and albeit not extensionally so (see previous \cref{rem:mu-congruence}).

% 
     We omit $\mu$-congruence for two main reasons: (a) we do not need it to obtain our ultimate undecidability results, rendering them strictly more general; and (b) the axiomatisation $\basicax$ remains purely equational, and so could be of independent interest.
\end{remark}


We shall soon present $\mu$-semirings equipped with an order, so as to axiomatise the \emph{leastness} of fixed points denoted by $\mu$, in particular satisfying the `strong' congruence principle mentioned in the example above.
Before that let us recall a class of models of $\basicax$ that is covered by our results:

\begin{definition}
[\cite{EsikLeiss05}]
    A \defname{Conway $\mu$-semiring} is a $\mu$-semiring further satisfying:
    \begin{enumerate}
    \item\label{item:fp-partial-unfolding} $\mu x e(f(x)) = e(\mu x f(e(x)))$
        \item\label{item:fp-diagonal} $\mu x \mu y e(x,y) = \mu x e(x,x)$
    \end{enumerate}
\end{definition}

Note that the fixed point axiom \eqref{item:fixed-point-axioms} is already a consequence of \eqref{item:fp-partial-unfolding} above, by setting $f(x)\df x$, and so Conway $\mu$-semirings are also models of $\basicax$.
Conway $\mu$-semirings were introduced in \cite{EsikLeiss05} as a minimal theory in which certain normal forms of context-free grammars may be formally obtained, in particular \emph{Greibach} normal form. 
We shall not develop its metatheory any further here, as our results about Conway $\mu$-semirings depend only on the fact that their equational theory is recursively enumerable, extends that of $\basicax$, and is modelled by the structure of languages $\Lang$ from \cref{language-model}.

\subsection{Equational systems and context-free grammars}
Before considering further extensions of $\basicax$, let us revisit the interpretation of $\mu$-expressions as context-free languages.
It is known that Conway $\mu$-semirings can `solve' context-free grammars (CFGs), but we shall recast (and improve) this result within $\basicax$ here.


An \defname{equational system} is a set $\mathcal E (\vec X) = \{X_i = f_i(\vec X)\}_{i<n}$.
 By recursively applying \emph{Beki\'c's Lemma} (\cite{Bekic84}) we can solve such systems in $\basicax$:

\begin{proposition}
\label{eq-systems-have-solns-in-basicax}
    Fix an equational system $\mathcal E(\vec X) = \{X_i = f_i(\vec X)\}_{i<n}$. 
    There are (closed) expressions $\vec e $ s.t.\ $\basicax \proves \mathcal E(\vec e)$. 
\end{proposition}
\begin{proof}
    By induction on $n$:
    \begin{itemize}
        \item When $n=0$ the statement is vacuously true.
        \item Fix an equational system $\{X_i = f_i(\vec X, X_n)\}_{i\leq n}$, where $\vec X = X_0, \dots, X_{n-1}$.
        Write $e_n'(\vec X) \df \mu X_n f_n (\vec X, X_n)$.
        Now, by inductive hypothesis, let $\vec e$ be solutions of $\{X_i = f_i(\vec X, e_n'(\vec X))\}_{i<n}$, i.e.\ such that $\basicax \proves e_i = f_i(\vec e, e_n'(\vec e))$, for all $i<n$.
        It suffices now to show that $e_n \df e_n'(\vec e) $ is a solution for $X_n$, i.e.\ $\basicax \proves e_n = f_n (\vec e, e_n)$, for which we reason as follows:
        \[
        \begin{array}{r@{\ = \ }ll}
             e_n & e_n' (\vec e) & \text{by definition of $e_n$} \\
                & \mu X_n f_n (\vec e, X_n) & \text{by definition of $e_n'$} \\
                & f_n (\vec e, \mu X_n f_n (\vec e, X_n)) & \text{by fixed point axiom \eqref{item:fixed-point-axioms}} \\
                & f_n (\vec e, e_n) & \text{by definition of $e_n$} \qedhere
        \end{array}
        \]
    \end{itemize}
\end{proof}

% Note that the result above does not rely on idempotency.
Recalling \cref{rem:mu-congruence}, note that solutions to an equational system are not necessarily unique, as $\basicax$ does not insist that $\mu$ computes \emph{least} fixed points.

Of course, we can use the above result to `solve' CFGs too.
% over $\Var$ a set of non-terminals $X,Y$ etc.\ generating languages over a finite alphabet $\Alphabet$ of letters $a,b,$ etc. We write $\alpha,\beta,$ etc.\ for strings of letters and non-terminals (i.e.\ elements of $(\Var \cup \Alphabet)^*$).

\begin{corollary}
\label{cfgs-have-solutions-in-basicaxs}
    Fix a CFG $\{X_i \Rightarrow f_{ij} (\vec X)\}_{i<n, j}$.
    There are $\mu$-expressions $\vec e = e_0, \dots, e_{n-1}$ s.t.\ $\basicax \proves e_i = \sum_j f_{ij} ({\vec e})$ for all $i<n$.
\end{corollary}

\subsection{Ordered semirings with fixed points}
Let us continue to recall certain classes of semirings over $\mu$-expressions from the literature.

    A semiring is \defname{idempotent} if it satisfies:
    \[
    \idem\ : \quad  e + e = e
    \]
We shall write $e\leq f \df e+f = f$,
typically in the presence of idempotency.
Note $\leq$, sometimes called the `natural order', defines a partial order under $\basicax + \idem$, and moreover one over which the semiring operations are monotone:

\begin{proposition}
\label{props-of-nat-order}
    The following hold in all idempotent semirings, and so in particular are consequences of $\basicax + \idem$:
% \begin{equation}
% \label{eq:po-props-needing-idem}
%     i\leq i 
%     \qquad 
%     e\leq f \implies e+i \leq f+i
% \end{equation}
%     \begin{equation}
%     \label{eq:po-props-not-needing-idem}
% \begin{array}{c}
%      % e\leq e \\
%      e \leq f \  \& \ f \leq e \implies e = f \\
%      e \leq f \ \& \  f \leq g \implies e \leq g
% \end{array}
% \qquad
% \begin{array}{c}
% % 0\leq e \\
% %     e_0 \leq f_0 \ \&\ e_1 \leq f_1 \implies e_0 + e_1 \leq f_0 + f_1 \\
% %     e_0 \leq f_0 \ \&\ e_1 \leq f_1 \implies e_0  e_1 \leq f_0  f_1
% e\leq f \implies eg \leq fg \\
% e \leq f \implies ge \leq gf
% \end{array}
% \end{equation}
\begin{equation}
    \label{eq:mon-po}
    \begin{array}{c}
     i\leq i \\
     e \leq f \  \& \ f \leq e \implies e = f \\
     e \leq f \ \& \  f \leq g \implies e \leq g
\end{array}
\qquad
\begin{array}{l}
e\leq f \implies e+i \leq f+i \\ 
e\leq f \implies eg \leq fg \\
e \leq f \implies ge \leq gf
\end{array}
\end{equation}
Moreover these properties hold in \emph{any} semiring, as long as $i$ is idempotent, and in particular are consequences of $\basicax + (i+i=i)$.
\end{proposition}
\begin{corollary}
\label{intro-elim-props-for-+}
    The following hold in all idempotent semirings, and so in particular are consequences of $\basicax + \idem$:
    \begin{equation}
        \label{eq:props-of-nat-ord}
    \end{equation}
    Moreover these properties hold in \emph{any} semiring, as long as $i$ is idempotent, and in particular are consequences of $\basicax + (i+i=i)$.
\end{corollary}


\begin{remark}
    [Ordered $\mu$-semirings]
    An \defname{ordered $\mu$-semiring} is a $\mu$-semiring equipped with a partial order satisfying the properties from \cref{eq:mon-po} above, as well as an ordered analogue of $\mu$-congruence, \defname{$\mu$-monotonicity}: $\dfrac{e(x)\leq f(x)}{\mu x e(x)\leq\mu x f(x)}$ (see \cite{EsikKuich12:mod-aut-theory,EsikLeiss02} for further details).
    
    Again, this does not necessarily hold in models of $\basicax+\idem$, even ones closed under (strong) $\mu$-congruence.
    For instance if $e(x) $ is $a+x$ and $f(x) $ is $a + b + x$, then $e(x)\leq f(x)$ is valid in all models of $\basicax+\idem$, but $\mu x e(x)$ needs only be greater than $a$ and $\mu x f(x) $ needs only be greater than $a+b$.
    Again tweaking the language model, it is consistent to interpret $\mu x e(x) $ as $\{a,b,c\}$ but $\mu x f(x)$ as $\{a,b\}$.
    Note that, since $e(x) \leq f(x) $ is valid in \emph{all} models of $\basicax + \idem + \mu$-congruence, it follows that $\mu$-monotonicity is not even admissible for this theory.
\end{remark}


    
In this work we shall only consider the `natural order' already defined, $e\leq f \df e+f = f$, before more generally structures without idempotency.
It is worth recalling at this point:
\begin{fact}
[Idempotent ideal]
\label{rem:idempotents-form-ideal}
The class of idempotent elements of a semiring always forms an ideal.
% The induced substructure is thus a model of $\basicax + \idem$, and in particular satisfies the properties of a monotone partial order in \cref{props-of-nat-order}.
\end{fact}


In the presence of an order it is of particular interest to examine structures where $\mu$ calculates a \emph{least} (pre)fixed point wrt.\ the order,
\[
\ind\ :\quad  
e(\mu x e(x)) \leq \mu x e(x)
\qquad \& \qquad
e(f)\leq f \implies \mu x e(x) \leq f
\]
yielding an induction principle.
Note that the axioms above are neither (in)equational, nor finitely many: there is an axiom for each choice of $f$ above.

\begin{definition}
        A \defname{Park $\mu$-semiring} is a $\mu$-semiring satisfying $\ind$. A \defname{Chomsky Algebra} is an idempotent Park $\mu$-semiring whose order is the natural order.\footnote{In fact the final clause here is redundant: the \emph{only} monotone partial order on an idempotent semiring is the natural order.}
\end{definition}

Chomsky Algebras were defined in \cite{GHK13,Leiss16} as \emph{algebraically complete} idempotent semirings, while Park $\mu$-semirings were defined axiomatically in \cite{EsikLeiss02,EsikLeiss05}. 
Both presentations are equivalent, when construing $\mu$-expressions as a naming convention for equational systems (see previous subsection), as was noticed already in \cite{GHK13}.
Note that the fixed point axiom \eqref{item:fixed-point-axioms} is already a consequence of $\ind$ (see \cite{EsikLeiss02,EsikLeiss05}), and so Chomsky Algebras are models of $\basicax+\idem$, and Park $\mu$-semirings are models of $\basicax$.


% \cite{EsikLeiss02,EsikLeiss05} considered more generally (not necessarily idempotent)
% Park $\mu$-semirings, i.e.\ $\mu$-semirings equipped with an order (not necessarily the natural order) over which all basic operations are monotone (see earlier Remark).

Again, we shall not recount the established metatheory of (idempotent) Park $\mu$-semirings or Chomsky Algebras: what is important for our results are that they are recursively enumerable extensions of $\basicax (+ \idem)$, that are still sound for $\Lang$.

% \input{decidability-mu-semirings}

\section{Effective inseparability of two problems on CFGs}
\label{sec:eff-insep}

Before presenting our main results on the undecidability of various theories of (idempotent) semirings with fixed points, we shall take a detour through the recursion theory of problems on CFGs.

\subsection{Two problems on CFGs}

For our arguments, we shall use {\em effective inseparability} of two algorithmic problems for context-free grammars (CFG). 
We shall consider only a special class of CFGs which we call \emph{productive} ones. Let $\Alphabet$ be the set of all terminal symbols and let $\Var$ be the set of all non-terminal ones.
As usual we shall assume among the non-terminals a designated \textbf{start} symbol $S \in \Var$.

\begin{definition}
[Productive CFGs]
A CFG is \textbf{productive} if, for every production rule $X \Rightarrow \alpha_1 \ldots \alpha_n$, some $\alpha_i$ is a terminal, i.e.\ some $\alpha_i \in \Alphabet$.
\end{definition}

A typical example of productive CFG is a Greibach grammar without $\varepsilon$-pro\-duc\-ti\-ons. In Greibach grammars, we always have $\alpha_1 \in \Alphabet$. Thus, by Greibach's normalisation theorem~\cite{Greibach1965}, one shows that any context-free language without the empty word is generated by a productive CFG. On the other hand, the empty word may never be generated by such a grammar.
In the sequel we shall restrict our attention to productive CFGs.

The first problem we are going to consider is the {\em totality} problem (i.e., the question whether a given productive CFG generates all non-empty words over $\Alphabet$). The second one is more specific. It could be called ``the grammar is total for obvious reasons,'' or {\em regularly total}. 
% and 
% the formal definition is given below. 
Formally:

%In what follows, we consider only context-free grammars of a specific form. Each grammar includes {\em three} designated symbols $S, Y, T$, where $S$ is the starting symbol, and the following fixed productions:
%\begin{align*}
%& S \Rightarrow YT \\
%& T \Rightarrow a T \mbox{ for each $a \in \Sigma$}\\
%& T \Rightarrow \varepsilon
%\end{align*}

\begin{definition}
[(Regular) totality]
\label{def-regtotal}
A productive CFG is \textbf{total}, if it generates (from $S$) all non-empty words over $\Alphabet$.
% \end{definition}

% \begin{definition}
% [Regular totality]

A productive CFG $G$ is \textbf{regularly total}, if it includes two designated non-terminals $Y,T \in \Var$ such that:
\begin{itemize}
    \item For each $a \in \Alphabet$, $G$ includes all productions of the form:
    \begin{align*}
& S \Rightarrow a Y T \\
& T \Rightarrow a \\
& T \Rightarrow a T 
\end{align*}
\item There is some $n>0$ such that:
\begin{enumerate}
    \item $S$ generates all nonempty words of length $\leq n+1$;
    \item $Y$ generates all words of length $n$.
\end{enumerate}
\end{itemize}

\end{definition}

It is not hard to see that
every regulary total productive CFG is total: 

\begin{proposition}
    [Regular totality implies totality]
    Every regularly total productive CFG is total.
\end{proposition}
\begin{proof}
    Let $G$ be a regularly total productive CFG and let $\vec a$ be a nonempty word:
    \begin{itemize}
        \item If $|\vec a| \leq n+1$ then it is generated from $S$ by definition of regular totality.
        \item If $|\vec a | > n+1$ then write $\vec a = a\vec b \vec c$, for some $a \in \Alphabet$, $|\vec b| = n$ and $|\vec c| > 0 $. 
        Then $S \Rightarrow a YT$ and $Y$ generates $\vec b$, by assumption, and $T$ generates $\vec c$ (since it trivially generates \emph{every} nonempty word). \qedhere
    \end{itemize}
\end{proof}

% Obviously, every regularly total productive CFG is total. Moreover, words of length at least $n+2$ are generated in a regular fashion, using the non-terminal $T$ which trivially generates all non-empty words. 
The algorithmic complexity of regular totality is dual to that of totality: while the latter is known to be $\Pi^0_1$-complete, the former belongs to $\Sigma^0_1$ (and, as we shall see below, is also complete in this class).

In what follows, we tacitly identify productive CFGs with their G\"odel numbers under a reasonable encoding. 
We write $\regtot$ for the set of (G\"odel numbers of) productive CFGs which are regularly total and $\nottot$ for the set of (G\"odel numbers of) productive CFGs which are not total.


\subsection{Effective inseperability by reduction}

We recall, following~\cite{Kuznetsov2021TOCL}, the notion of effective inseparability and the results which allows it to be used for proving $\Sigma^0_1$-completeness of logical theories. 

Fix some acceptable enumeration $\{ \prf n\}_{n\in \N}$ of the one-input partial recursive functions on $\N$, 
and let us write $\res n$ for ``the $n$\textsuperscript{th} r.e.\ set,'' i.e.\ $\dom { \prf n}$, the domain of $ \prf n$.

\begin{definition}
$A,B\subseteq \N$ are \textbf{effectively inseparable} if:
\begin{itemize}
    \item $A \cap B = \emptyset$; and,
    \item there is a partial recursive function $h(x,y)$ s.t., whenever $\res u \supseteq A$ and $\res v \supseteq B$ with $\res u \cap \res v = \emptyset$, we have $h(u,v) \downarrow\,  \notin \res u \cup \res v$. 
\end{itemize}
We shall call such $h(x,y)$ a \textbf{productive function} for $A,B$.
\end{definition}


Given $A,B \subseteq \N$, a set $C\subseteq \N $ \textbf{separates} $A$ from $B$ (or is a \textbf{separator}) if $C\supseteq A$ but nonetheless $C\cap B = \emptyset$. 
In the Definition above, informally, $h$ uniformly falsifies any attempt to separate $A$ from $B$ by a decidable set.
Intuitively, we can see $\lambda v\,  h(u,v)$ here as a `proof' that any such $\res u$ is not decidable: it is a recursive function that computes counterexamples to the statement that $\res u$ has an r.e.\ complement.

The following result is a corollary of a theorem by Myhill~\cite{Myhill1955}, as shown in~\cite[Exercise 11-14]{Rogers1987}:
\begin{theorem}
\label{separator-of-eff-insep-sets-is-S01-complete}
If $A,B \subseteq \N $ are effectively inseparable r.e.\ sets and $C$ is an r.e.\ set separating them then $C$ is $\Sigma^0_1$-complete.
\end{theorem}

Note in the Theorem above, and generally in the sequel, we shall not be precise about whether $C$ separates $A$ from $B$ or vice-versa, WLoG by symmetry of the statement.


Effective inseparability can be propagated by a (very weak) kind of reduction.
Given $A,B\subseteq \Nat$ with $A\cap B = \emptyset$ and $A', B'\subseteq \Nat$ with $A' \cap B' = \emptyset$, an \textbf{e.i.-reduction} from $(A',B')$ to $(A,B)$ is a total computable function $f$ such that: 
\begin{enumerate}
% \item $x \in A \implies f(x) \in A'$; 
\item\label{item:f(A)-in-A'} $f(A) \subseteq A'$
% \item $x \in B \implies f(x) \in B'$.
\item\label{item:f(B)-in-B'} $f(B) \subseteq B'$
\end{enumerate}
Such maps are indeed appropriate for reducing questions of effective inseperability:

\begin{proposition}\label{prop-reduction}
Suppose $A,B \subseteq \N$ are effectively inseparable, and 
 $A',B' \subseteq \N$ with $A' \cap B' = \emptyset$.  
 % and $f$ is a total computable function s.t.,
% \begin{enumerate}
% % \item $x \in A \implies f(x) \in A'$; 
% \item\label{item:f(A)-in-A'} $f(A) \subseteq A'$,
% % \item $x \in B \implies f(x) \in B'$.
% \item\label{item:f(B)-in-B'} $f(B) \subseteq B'$,
% \end{enumerate}
If $f$ is an e.i.-reduction from $(A',B')$ to $(A,B)$ 
then $A'$ and $B'$ are also effectively inseparable.
\end{proposition}

\begin{proof}
The preimage of an r.e.\ set under a total computable mapping is itself r.e.
Indeed we have $\inv f (\res n) = \dom {\prf n \circ f}$, and we can even effectively compute from $n$ the index of the r.e.\ function $\prf n\circ f $, say by a total recursive function $g$, so that always $\inv f(\res n) = \res {g(n)}$.
% Indeed, if $W_n$ is the domain of a partial computable function, then $f^{-1}(W_n)$ is the domain of the composition of $f$ with this function. Moreover, for any total computable function $f$ there exists a total computable function $g$ such that $f^{-1}(W_n) = W_{g(n)}$. 

Now, let $h$ be a productive function for $A,B$ and let us define a productive function $h'$ for $A',B'$ as follows: \[h'(u,v) = f(h(g(u), g(v)))\] 

Clearly $h'$ is a partial recursive function.
To show it is productive for $A',B'$, suppose
  $\res u \supseteq A'$, $\res v \supseteq B'$ with
  $\res u \cap \res v = \emptyset$ and let us show that $h'(u,v) \notin \res u \cup \res v$. 
First we have:
\[
\begin{array}{rcll}
     \res{g(u)} & = & \inv f (\res u ) & \text{by definition of $g$} \\
        & \supseteq & \inv f (A') & \text{since $\res u \supseteq A'$} \\
        & \supseteq & A & \text{since $f(A) \subseteq A'$} \\
    \noalign{\medskip}
    \res{g(v)} & = & \inv f (\res v ) & \text{by definition of $g$} \\
        & \supseteq & \inv f (B') & \text{since $\res v \supseteq B'$} \\
        & \supseteq & B & \text{since $f(B) \subseteq B'$} \\
    \noalign{\medskip}
    \res{g(u)} \cap \res{g(v)} & = & \inv f (\res u) \cap \inv f (\res v) & \text{by definition of $g$} \\
        & = & \inv f (\res u \cap \res v ) & \text{since preimages preserve intersections}\\
        & = & \inv f (\emptyset) & \text{since $\res u \cap \res v = \emptyset$} \\
        & = & \emptyset
\end{array}
\]

% By \cref{item:f(A)-in-A',item:f(B)-in-B'} we have $f^{-1}(A') \supseteq A$ and $f^{-1}(B') \supseteq B$.

%   Thus, for $W_{g(u)} = f^{-1}(W_u)$ and $W_{g(v)} = f^{-1}(W_v)$ we have the following:
% \begin{align*}
% & W_{g(u)} \supseteq f^{-1}(A') \supseteq A; \\
% & W_{g(v)} \supseteq f^{-1}(B') \supseteq B; \\
% & W_{g(u)} \cap W_{g(v)} = \varnothing.
% \end{align*}

From here we have:
\[
\begin{array}{rcll}
     h(g(u),g(v)) & \notin & \res{g(u)} \cup \res{g(v)} & \text{since $h$ is productive for $A,B$} \\
        & \notin & \inv f (\res u) \cup \inv f (\res v) & \text{by definition of $g$} \\
        & \notin & \inv f (\res u \cup \res v) & \text{since preimages preserve unions}
\end{array}
\]
% 
% Now, since $h$ witnesses effective inseparability of $A$ and $B$, we have $h(g(u),g(v)) \notin W_{g(u)} \cup W_{g(v)} = 
% f^{-1}(W_u \cup W_v)$. 
Hence $h'(u,v) = f(h(g(u), g(v)) \notin  \res u \cup \res v$, as required.
\end{proof}



\subsection{Effective inseperability of $\regtot$ and $\nottot $}

The main result of this section is the following:

\begin{theorem}
\label{regtot-nottot-eff-insep}
The sets $\regtot$ and $\nottot$ are effectively inseparable.
\end{theorem}
This subsection is devoted to proving this result.
We proceed via a reduction from a more well-known pair of effectively inseparable problems: halting and looping for Turing machines. 
% 
Throughout this section we shall work with one-input deterministic Turing machines that include, besides the initial and final state, also a designated {\em looping} state $q_c$ with rules that force the machine to remain at this state once it is reached: $(q_c, a) \Rightarrow (q_c, a, N)$ for any letter $a$ of the tape alphabet.

% Let $\TM$ and $x$ denote an arbitrary Turing machine and its input word, respectively. 
We say that a machine $\TM$ loops on an input $x$ if $\TM$ reaches $q_c$ when running on $x$. Looping is a very specific case of non-halting. Let 
\begin{align*}
& \loops = \{ (\TM,x) \mid \mbox{$\TM$ loops on input $x$} \}, \\
& \halts = \{ (\TM,x) \mid \mbox{$\TM$ halts on input $x$} \}
\end{align*}
Again, we construe both these sets as subsets of $\Nat$ by implicitly identifying pairs $(\TM,x)$ with their codes.
The following fact is folklore (see, e.g., \cite[Exercise 7-55d]{Rogers1987}\cite[Proposition 3]{Kuznetsov2021TOCL}):

\begin{proposition}\label{prop-insepCH}
The sets $\loops$ and $\halts$ are effectively inseparable.
\end{proposition}

Now we can prove the effective inseparability of $\regtot$ and $\nottot$ by reduction, in the sense of Proposition~\ref{prop-reduction}, to the effective inseparability of $\loops$ and $\halts$.


Let us start by briefly recalling some ideas from the standard proof of $\Pi^0_1$-hardness for the totality problem for CFGs, as presented in Kozen's textbook~\cite{Kozen1997}. 
The {\em halting protocol} of a Turing machine $\TM$ on a given input $x$ (if it exists) is the word $\# \kappa_0 \# \kappa_1 \# \ldots \# \kappa_n \#$, where $\kappa_0$ is the initial configuration $q_0 x$, each $\kappa_{i+1}$ is the configuration which is the immediate successor of $\kappa_i$, and $\kappa_n$ is a final configuration. Let $\Alphabet$ be the alphabet in which configurations are written down (which includes the internal alphabet of $\TM$, its set of states, and the additional symbol $\#$).
We may readily define a (computable) {reducing function}, mapping each pair $(\TM,x)$ to a CFG which generates all non-empty words except for the halting protocol (if it exists). Thus, this grammar is total if and only if $\TM$ does not halt on $x$. 



We shall modify this construction a bit, following~\cite{Kuznetsov2021TOCL}, in order to take care of looping. 
All non-empty words which are not the halting protocol must belong to one of the following three classes:%\anupam{(I) and (II) are not necessarily disjoint, right? If so we could mention/emphasise this.}
\begin{enumerate}[(I)]
\item\label{not-prefix-haltprot} words of length greater than 1, beginning with $\#$, which cannot be even a prefix of the halting protocol: this includes words not beginning with $\# \kappa_0 \#$; words where there is not a correct configuration between two $\#$'s; words with subwords $\# \kappa \# \kappa' \#$ where $\kappa'$ is not the immediate successor of $\kappa$; and, most importantly, all words including $q_c$;
\item\label{incomplete-prot} words of length greater than 1, beginning with $\#$, and not ending with $\# \kappa_F \#$ where $\kappa_F$ is a final configuration (such a word could be an incomplete protocol---a prefix of the halting protocol);
\item\label{begins-incorrectly} words beginning with a letter different from $\#$, and the one-letter word $\#$.
\end{enumerate}
Note that these classes are not necessarily disjoint: a word can belong to more than one of them.

Write (I)$'$ and (II)$'$ for the sets of words from (I) and (II), respectively, with the first symbol $\#$ removed.
The standard construction shows that both (I)$'$ and (II)$'$ are context-free. Denote the grammars generating (I)$'$ and (II)$'$ by $G_1$ and $G_2$, respectively; we suppose that both grammars are productive (e.g., being Greibach grammars). Here it is important that (I)$'$ and (II)$'$ are $\varepsilon$-free. Moreover, $G_1$ and $G_2$ are computable from $(\TM,x)$~\cite[Section~2.1]{Kuznetsov2021TOCL}. 

We assume $G_1$ and $G_2$ have disjoint sets of non-terminals, writing $Y$ and $Z$ for their start symbols, respectively. 
Now define the grammar $G_{\TM,x}$ obtained from the union of $G_1$ and $G_2$ along with the following additional productions:
\begin{align*}
& S \Rightarrow aYT && \mbox{ for each $a \in \Alphabet$} \\
& T \Rightarrow a T && \mbox{ for each $a \in \Alphabet$} \\
& T \Rightarrow a && \mbox{ for each $a \in \Alphabet$} \\
& S \Rightarrow \# Y \\
& S \Rightarrow \# Z \\
& S \Rightarrow a T && \mbox{ for each $a \in \Alphabet - \{ \# \}$} \\
& S \Rightarrow a && \mbox{ for each $a \in \Alphabet$}
\end{align*}
Here $S$  and $T$ are fresh non-terminals, and $S$ is the new start symbol of $G_{\TM,x}$.
Notice that the first three productions are exactly as in \cref{def-regtotal} for regular totality.

\begin{proposition}
\label{grammar-generates-all-but-haltprot}
    $G_{\TM,x}$ generates all nonempty words over $\Alphabet$, except the halting protocol of $\TM$ on $x$ (if it exists).
\end{proposition}
\begin{proof}
For words from (I) or (II) we start by firing $S \Rightarrow \# Y$ and $S \Rightarrow \# T$, respectively, and continue as in $G_1$ or $G_2$, respectively. 
Words from (III) are generated by the final two rules. Note that this analysis accounts for \emph{all} words generated by runs beginning with the final four rules from the start symbol $S$.

It remains to verify that the halting protocol cannot be generated by runs beginning with the first rule $S \Rightarrow a YT$.
However, any such run generates only words not beginning with $\#$ and words with a prefix from (I); by definition, no such word could be the halting protocol.    
\end{proof}


We can now define an appropriate \emph{reducing function}:
\begin{equation}
\label{eq:ei-reduction-rnt-to-ch}
    f \colon (\TM,x) \mapsto G_{\TM,x}
\end{equation}
Of course, this function is total and computable.
Moreover it allows us to deduce effective inseparability of $\regtot $ and $\nottot$ from that of $\loops$ and $\halts$:



\begin{proposition}
\label{regtot-nottot-reduces-to-loops-halts}
    The function $f$ in \eqref{eq:ei-reduction-rnt-to-ch} is an e.i.-reduction from $(\regtot, \nottot)$ to $(\loops,\halts)$.
\end{proposition}
\begin{proof}
Clearly $\regtot \cap \nottot = \emptyset$. 
So we need to show the following two properties:
\begin{itemize}
\item if $(\TM,x) \in \mathbf{C}$, then $f((\TM,x)) = G_{\TM,x} \in \mathbf{R}$;
\item if $(\TM,x) \in \mathbf{H}$, then $f((\TM,x)) = G_{\TM,x} \in \nottot$.
\end{itemize}

The second property follows immediately from \cref{grammar-generates-all-but-haltprot}, since the halting protocol cannot be generated by $G_{\TM,x}$. The first property is more involved.

Let $\TM$ loop on $x$ and let $n+1$ be the length of an (incomplete) protocol of $\TM$ on $x$ which reaches a configuration where $\TM$ is in state $q_c$. We shall show that $G_{\TM,x}$ is regularly total, according to the given $n$.
That is, we will show that:
\begin{enumerate}
    \item $S$ generates all nonempty words of length $\leq n+1$.
    \item $Y$ generates all words of length $n$.
\end{enumerate}

For (1), since $\TM$ does not halt on $x$, there is no halting protocol, and $S$ generates all nonempty words, by \cref{grammar-generates-all-but-haltprot}. In particular, $S$ generates all words of length $\leq n + 1$.

For (2), we need to show that any word $\vec a$ of length $n$ belongs to (I)$'$.
For this, even if $\# \vec a$ is a prefix of the (infinite) protocol of $\TM$ on $x$, it includes, by choice of $n$, the letter $q_c$, and so $\# \vec a \notin $(I).
% 
% Indeed, take such a word and prepend it with $\#$. Even if the result is a prefix of the (infinite) protocol of $\TM$ on $x$, it includes, by choice of $n$, the letter $q_c$. Otherwise, there are some local conflicts which also put the word into (1). Thus, the original word is always generated by $Y$.
% 
%
%belongs to language (1): even if it is an incomplete protocol, it includes, by choice of $n$, the letter $q_c$. Thus, $Y$ generates all words of length $n$. 
% 
% By Definition~\ref{def-regtotal}, if $\TM$ loops on $x$, then $G_{\TM,x}$ is regularly total, i.e., $G_{\TM,x} \in \mathbf{R}$. 
\end{proof}

Now we can finally conclude the main result of this section:
\begin{proof}
    [Proof of \cref{regtot-nottot-eff-insep}]
    Follows immediately from \cref{prop-reduction,prop-insepCH,regtot-nottot-reduces-to-loops-halts}.
\end{proof}

% Now, by Propositions~\ref{prop-reduction} and~\ref{prop-insepCH} we conclude that $\mathbf{R}$ and $\overline{\mathbf{T}}$ are effectively inseparable.



\section{Undecidability of (extensions of) $\basicax + \idem$ }
\label{sec:undec-with-idem}

This section is devoted to a proof of our first main result:

\begin{theorem}
\label{thm:undecidability-with-idem}
    Any r.e.\ extension of $\basicax + \idem$ that is sound for $\Lang$ is $\Sigma^0_1$-complete.
\end{theorem}

We shall actually strengthen this result in the next section to avoid the need for idempotency (and even identity), by a sort of reduction, but we present this weaker version to illustrate the main ideas in a simpler setting.


\subsection{Some intermediate lemmata}
% For convenience, we shall identify non-terminals of a CFG with the corresponding expressions solving the CFG in $\basicax$, cf.~\cref{prop:cfgs-have-solutions-in-basicaxs}.

To prove \cref{thm:undecidability-with-idem} we first establish a couple intermediate results reasoning about (productive) CFGs in $\basicax (+ \idem)$.
As we shall later consider settings without idempotency and identity, we shall be careful to avoid relying on the multiplicative identity $1$, and shall also be sensitive to the use of idempotency.


\begin{lemma}
\label{lem:membership}
    Fix a CFG $G $ and write $ e_X$ for (any) corresponding solutions in $\basicax$ of each non-terminal $X$ of $G$, cf.~\cref{cfgs-have-solutions-in-basicaxs}. 
    If $X$ generates a word $\vec a$ then $\basicax + \idem \proves \vec a \leq e_X$.
\end{lemma}
\begin{proof}
    By induction on the derivation tree of $\vec a $ from $X$ in $G$.
    Working in $\basicax+\idem$:
    \begin{itemize}
        \item If $X \Rightarrow \vec a$, then we have $\vec a + g = e_X$ for some $g$ by \cref{cfgs-have-solutions-in-basicaxs}, and so $\vec a \leq e_X$ by idempotency of $\vec a$ and \cref{props-of-nat-order}.
        \item If $X \Rightarrow \vec a_0 X_1 \vec a_1 \cdots \vec a_{n-1}  X_n \vec a_n$, s.t.\ $\vec a = \vec a_0 \vec b_1 \vec a_1 \cdots \vec b_n \vec a_n$, then each $\vec b_i $ is generated by $X_i$ by smaller derivation trees.
        Thus by inductive hypothesis we have $\vec b_i \leq e_{X_i}$ for all $i$.
        Finally by monotonicity of products wrt $\leq$, we have indeed $\vec a \leq e_X$.\qedhere
    \end{itemize}
\end{proof}

Let us write $\top \df \mu x \left(\sum\limits_{a \in \Alphabet} a + \sum\limits_{a\in \Alphabet} ax \right)$, which trivially computes the `total' language of all nonempty words.

\begin{lemma}
\label{lem:unfolding-top}
    For $n>0$, we have $\basicax \proves \top = \sum\limits_{|\vec a|=1}^{n}\vec a + \sum\limits_{|\vec a|= n} \vec a \top$.
\end{lemma}

\begin{proof}
    By induction on $n>0$:
    \begin{itemize}
        \item When $n=1$, then the statement follows by just an unfolding of the fixed point $\top$, axiom \eqref{item:fixed-point-axioms}.
        \item For the inductive step we have:
        \[
        \begin{array}{r@{\ = \ }ll}
             \top & \sum\limits_{|\vec a|=1}^{n}\vec a + \sum\limits_{|\vec a|= n} \vec a \top & \text{by inductive hypothesis} \\
                & \sum\limits_{|\vec a|=1}^{n}\vec a + \sum\limits_{|\vec a|= n} \vec a \left(\sum\limits_{a \in \Alphabet} a + \sum\limits_{a\in \Alphabet} a\top \right) & \text{by unfolding $\top$, axiom \eqref{item:fixed-point-axioms}} \\
                & \sum\limits_{|\vec a|=1}^{n}\vec a + \sum\limits_{|\vec a|= n} \vec a \sum\limits_{a \in \Alphabet} a + \sum\limits_{|\vec a|= n} \vec a\sum\limits_{a\in \Alphabet} a\top & \text{by distributivity} \\
                & \sum\limits_{|\vec a|=1}^{n}\vec a + \sum\limits_{|\vec a|= n+1} \vec a + \sum\limits_{|\vec a|= n+1} \vec a\top & \text{by distributivity} \\
                & \sum\limits_{|\vec a|=1}^{n+1}\vec a + \sum\limits_{|\vec a|= n+1} \vec a\top & \text{by rebracketing sums} \qedhere
        \end{array}
        \]
    \end{itemize}
\end{proof}

\subsection{Small proofs of regular totality}

Note that $\top =\mu x \left(\sum\limits_{a \in \Alphabet} a + \sum\limits_{a\in \Alphabet} ax \right)$ is a solution of the non-terminal $T$ in any regularly total grammar.
In what follows, whenever we fix solutions of a regularly total grammar, we assume the solution for $T$ is $\top$.
Again, we shall be sensitive to the use of idempotency, making clear any such instances for our strengthenings in the next section.

\begin{proposition}
\label{reg-tot-provable-with-idem}
    Fix a productive CFG $G \in \mathbf R$ and let $e_X$ be solutions of each if its non-terminals $X$ in $\basicax$, cf.~\cref{cfgs-have-solutions-in-basicaxs}, with $e_T = \top$. 
    Then $\basicax + \idem \proves \top \leq e_S$.
\end{proposition}
\begin{proof}
Since $G\in \mathbf R$, let $n>0$ be such that:
\begin{enumerate}
    \item\label{item:S-gens-all-short-words} $S$ generates all nonempty words of length $\leq  n+1$.
    \item\label{item:Y-gens-bar} $Y$ generates all words of length $n$.
    \item\label{item:G-reaches-loop} $G$ has a production $S \to aYT$ for each $a \in \Alphabet$.
\end{enumerate}
    Let us work in $\basicax + \idem$. 
    By \eqref{item:Y-gens-bar} above and \cref{lem:membership} we have
    $ \vec a \leq e_Y$ whenever $|\vec a| = n$.
    By \purple{idempotency of $e_Y$} and \cref{intro-elim-props-for-+} we have $\sum\limits_{|\vec a| = n}\vec a \leq e_Y $.
    Thus by monotonicity of $\cdot$ and distributivity:
    $$\sum\limits_{|\vec a| = n}\vec a \top \leq e_Y\top$$
    Now by \eqref{item:G-reaches-loop} and \cref{cfgs-have-solutions-in-basicaxs} we have $g + \sum\limits_{a \in \Alphabet}ae_Y\top = e_S$ for some $g$ (and since $e_T$ is $\top$).
    So by \purple{idempotency of $\sum\limits_{a \in \Alphabet}ae_Y\top$} we have:
    $$\sum\limits_{a \in \Alphabet} ae_Y\top \leq e_S$$
    Putting the two displays above together we have $\sum\limits_{a \in \Alphabet} a\sum\limits_{|\vec a| = n}\vec a \top \leq e_S$ by monotonicity, and so by distributivity:
    \begin{equation}
    \label{eq:long-words-reach-S}
        \sum\limits_{|\vec a| = n+1}\vec a \top \leq e_S
    \end{equation}
    
    Now by \eqref{item:S-gens-all-short-words} and again by \cref{lem:membership} we have $\vec a \leq e_S$ whenever $1 \leq |\vec a|\leq n+1$. 
    Again by \purple{idempotency of $e_S$} and monotonicity of $+$ we have 
    $\sum\limits_{|\vec a|=1}^{n+1} \vec a \leq e_S$.
    Putting this together with \eqref{eq:long-words-reach-S} above, we have \purple{by idempotency of $e_S$} and \cref{intro-elim-props-for-+}:
    $$\sum\limits_{|\vec a|=1}^{n+1} \vec a  + \sum\limits_{|\vec a| = n+1}\vec a \top \leq e_S$$
    From here we conclude by \cref{lem:unfolding-top} and transitivity of $\leq$, \cref{props-of-nat-order}.
\end{proof}

\subsection{Putting it together}
We are now ready to prove the main result of this section:
% \begin{theorem}
%     Any r.e.\ extension of $\basicax + \idem$ that is sound for all $\mu$-continuous Chomsky Algebras (even the equational theory of $\Lang$) is $\Sigma^0_1$-complete.
% \end{theorem}

\begin{proof}
    [Proof of \cref{thm:undecidability-with-idem}]
    Let us assume that all CFGs have start symbol $S$ and write $e^G_X$ for the solution of the nonterminal $X$ in $G$ obtained by \cref{cfgs-have-solutions-in-basicaxs}.
    Let $\mathsf U \supseteq \basicax+\idem$ be r.e.\ with $\Lang \models \mathsf U$.  
    Writing $\total {\mathsf U} \df \{G \text{ productive} : \mathsf U \proves \top \leq e^G_S \}$,
    we have:
    \begin{itemize}
        \item $\total {\mathsf U}  \cap \nottot = \emptyset$, by assumption that $\Lang \models \mathsf U$.
        \item $\mathbf R \subseteq \total {\mathsf U}  $, since $\basicax+\idem \subseteq \mathsf U$ and by \cref{reg-tot-provable-with-idem}.
    \end{itemize}
    Thus by \cref{separator-of-eff-insep-sets-is-S01-complete,regtot-nottot-eff-insep} we have that $\total {\mathsf U}$ is $\Sigma^0_1$-complete.
    Since there is an obvious computable reduction from $\total {\mathsf U}$ to $\mathsf U$, we have that $\mathsf U$ is $\Sigma^0_1$-complete too, as required.
\end{proof}


Since all theories we have introduced are r.e.\ and modelled by $\Lang$, we thus immediately have:
\begin{corollary}
    The equational theories of idempotent Conway $\mu$-semirings, idempotent Park $\mu$-semirings and Chomsky algebras are undecidable, and moreover $\Sigma^0_1$-complete.
\end{corollary}

\section{The case without idempotency}
\label{sec:undec-without-idem}

In this section we shall strengthen the main result of the previous section to cover theories of semirings that are not necessarily idempotent.

% \subsection{Productive grammars}
% Call a CFG \emph{productive} if, for every transition $X \to \vec \alpha$, with each $\alpha_i \in \Alphabet \cup \Var$, there is some $\alpha_i \in \Alphabet$.

% For example, any Greibach grammar without $\epsilon$-productions is productive.
% It is thus not difficult to show that all $\epsilon$-free CFLs are generated by productive (even $\epsilon$-free Greibach) grammars.


% \begin{definition}
%     Write $\tot^+$ for the class of productive CFGs that generate all nonempty words.
    
%     Write $\regtot^+$ for the class of productive CFGs with start symbol $S$ and including nonterminals $Y,T$ s.t.\ there is some $n>0$ with:
%     \begin{itemize}
%         \item $S$ generates all nonempty words of length $\leq n$; 
%         \item there are productions $S \to aY$ and $S \to aYT$ for each $a \in \Alphabet$;
%         \item $Y$ generates every word of length $n$.
%         \item there are productions $T \to a$ and $T \to aT$ for each $a \in \Alphabet$.
%     \end{itemize}
% \end{definition}

% It is not hard to see that $\regtot^+ \subseteq \tot^+$.
% Moreover we want:

% \begin{theorem}
%     $\regtot^+$ and $\overline \tot^+$ are effectively inseparable.
%     \todo[inline]{to prove!}
% \end{theorem}

\subsection{Identity-free $\mu$-semirings and productive expressions}

Write $\basicaxidfree$ for the axiomatisation defined just like $\basicax$ but without any of the axioms mentioning $1$ (and without $1$ in the language).

The \defname{productive} expressions (in the language of $\basicaxidfree$), written $p,q$ etc., are generated by,
\begin{equation}
    \label{eq:productive-expressions}
    p,q,\dots \quad ::= \quad 0 \quad \mid \quad a 
% \quad \mid \quad pX \quad \mid \quad Xp 
\quad \mid \quad p+q 
% \quad \mid \quad pq 
\quad \mid \quad pe \quad \mid \quad ep
\quad \mid \quad \mu X p
\end{equation}
where $e$ ranges over (possibly open, not necessarily productive) $1$-free expressions.

By inspection of the proof of \cref{eq-systems-have-solns-in-basicax}, we can give a more refined result than \cref{cfgs-have-solutions-in-basicaxs}:

\begin{proposition}
\label{prod-cfg-has-prod-solns}
    A productive CFG has productive solutions in $\basicaxidfree$.
    I.e., given a productive CFG $\{X_i \Rightarrow f_{ij}(\vec X)\}_{i<n,j}$, there are productive expressions $\vec e = e_0, \dots, e_{n-1}$ s.t.\ $\basicaxidfree \proves e_i = \sum_j f_{ij}(\vec e)$ for all $i<n$.
\end{proposition}
In particular note that the RHS of a production in a productive CFG is a productive expression, whence the same proof of \cref{eq-systems-have-solns-in-basicax} applies for the statement above.
Note also that that proof does not rely on any properties of the multiplicative identity $1$ (which in any case plays no role in productive CFGs).

\subsection{An idempotent interpretation of productive expressions}

Temporarily write $e':= \mu x (e + \mu y (x + y))$. 
It is not hard to see that, over $\basicaxidfree$, $e'$ is an idempotent:
\[
\begin{array}{r@{\ = \ }ll}
      e'  & e + \mu y (e' + y) & \text{by unfolding $e'$, by axiom \eqref{item:fixed-point-axioms}} \\
        & e + e' + \mu y (e' + y) & \text{by unfolding $\mu y (e' + y)$, by axiom \eqref{item:fixed-point-axioms}} \\
        & (e + \mu y (e' + y) ) + e' & \text{by associativity and commutativity of $+$}\\
        & e' + e' & \text{by the first line}
\end{array}
\]

In the presence of an order $e'$ is also an upper bound of $e$, and in the presence of induction $e'$ is in fact the \emph{least} idempotent upper bound of $e$.
In particular in $\Lang $ (and any other idempotent model) we have simply $e'=e$:
\begin{itemize}
    \item On one hand, we have $\lang {e'} = \lang{e + \mu y (e'+y)}$, and so $\lang {e'} \supseteq \lang e$.
    \item On the other hand, we have that $e$ solves the matrix of $e'$:
    \begin{itemize}
        \item Clearly $\lang {\mu y (e+y)} = e$: as above we have $\lang {\mu y (e+y)} = \lang{e + \mu y (e+y)} \supseteq \lang e$; we also have that $e$ solves the matrix of $\mu y (e+y)$ in $\Lang$, i.e., $\lang e = \lang {e+e}$, thanks to idempotency.
        \item Thus we have $e = e + e = e + \mu y (e+y)$ in $\Lang$, again by idempotency.
    \end{itemize}
    Thus indeed $\lang {e'}\subseteq \lang e$, as $e'$ must be a least solution.
\end{itemize}

\begin{definition}
    [Idempotent interpretation]
    For $1$-free expressions $e$, define $\id e $ as follows:
\begin{itemize}
    \item $\id a \df a' = \mu x (a + \mu y (x+y))$
    \item $\id \cdot $ commutes with everything else. I.e.,
    \[
    \begin{array}{r@{\ \df\ }l}
         \id 0 & 0 \\
         \id x & x \\
         \id {(e+f)} & \id e + \id f \\
         \id {(ef)} & \id e \id f \\
         \id{(\mu x e)} & \mu x \id e
    \end{array}
    \]
%     \begin{itemize}
%         \item $\id 0 \df 0$
%         % \item $\id 1 \df 1$
%         \item $\id x \df x$
%         \item $\id {(e+f)} \df \id e + \id f$
%         \item $\id {(ef)} \df \id e \id f$
%         \item $\id {(\mu x e)} \df \mu x \id e$
%     \end{itemize}
\end{itemize}
\end{definition}

As previously argued it is not hard to see that:
\begin{proposition}
\label{idem-trans-equiv-in-lang}
    $\lang {\id e} = \lang {e}$.
\end{proposition}
\begin{proof}
[Proof sketch]
    By induction on the structure of $e$. In fact, we must strengthen the statement to $\lang {\id e(\vec A)} = \lang {e(\vec A)}$ for all languages $\vec A$.
    The case when $e$ is some $a \in \Alphabet$ follows from the argument above.
    Every other case is immediate, with the $\mu$-case following by the strong $\mu$-congruence principle in $\Lang$: $\forall x\,    e(x) =f(x) \limp \mu x e(x) = \mu x f(x)$.
\end{proof}

Moreover, $\id \cdot$ induces a bona fide interpretation of identity-free $\mu$-semirings:

\begin{proposition}
\label{idem-trans-is-an-interp}
    $\basicaxidfree\proves e=f \implies \basicaxidfree \proves \id e = \id f$.
\end{proposition}
\noindent
To see this, simply replace every expression $e $ in a $\basicaxidfree$ proof by $\id e$. It remains to verify that the axioms and rules remain correct, which is routine.




Finally, thanks to the fact that idempotent elements of a semiring form an ideal, cf.~\cref{rem:idempotents-form-ideal}, we have idempotency of the image of $\id \cdot $ on productive expressions provably in $\basicaxidfree$:

\begin{lemma}
\label{idem-trans-is-idem-on-prod-exprs}
    If $e(\vec x)$ is productive then $\basicaxidfree \proves \id e(\vec f) + \id e(\vec f) = \id e(\vec f)$ for all $\vec f$.
\end{lemma}
\begin{proof}
    By induction on the structure of $e(\vec x)$, according to the grammar of productive expressions in \eqref{eq:productive-expressions}.
    Working in $\basicaxidfree$ we have:
    \begin{itemize}
        \item If $e(\vec x)$ is $0$ then: 
        \[
        \begin{array}{r@{\ = \ }ll}
             \id 0 + \id 0 & 0 + \id 0 & \text{by definition of $\id 0$} \\
                & \id 0 & \text{since $0$ is an additive unit}
        \end{array}
        \]
        \item If $e(\vec x)$ is some $a\in \Alphabet$ then $\id a + \id a = \id a$ by the argument at the beginning of this subsection.
        \item If $e(\vec x)$ is $p(\vec x) +  q(\vec x) $ then:
        \[
        \begin{array}{rll}
        % & \id{(p(\vec f)+q(\vec f))} + \id{(p(\vec f)+q(\vec f))} \\
        % = 
        & (\id p(\vec f) + \id q(\vec f)) + (\id p(\vec f) + \id q(\vec f)) & 
        % \text{by definition of $\id \cdot$}
        \\
          =& (\id p(\vec f) + \id p(\vec f)) + (\id q(\vec f) + \id q(\vec f)) & \text{by associativity and commutativity of $+$} \\
            =& \id p(\vec f) + \id q(\vec f) & \text{by inductive hypothesis for $p(\vec x)$ and $q(\vec x)$} 
            % \\
            % =& \id{(p(\vec f)+q(\vec f))} & \text{by definition of $\id \cdot$}
        \end{array}
        \]
        \item If $e(\vec x)$ is $p(\vec x)g(\vec x)$ then:
        \[
        \begin{array}{rll}
        % \id{(p(\vec f)g(\vec f))} + \id{(p(\vec f)g(\vec f))} 
        &  \id p(\vec f)\id g(\vec f) + \id p(\vec f) \id g(\vec f) & 
        % \text{by definition of $\id\cdot$} 
        \\
          =  & (\id p(\vec f) + \id p(\vec f))\id g(\vec f) & \text{by distributivity} \\
           = & \id p(\vec f) \id g(\vec f) & \text{by inductive hypothesis for $p(\vec x)$} 
           % \\
           % = & \id {(p(\vec f)g(\vec f))} & \text{by definition of $\id \cdot$}
        \end{array}
        \]
        \item If $e(\vec x)$ is $g(\vec x)p(\vec x)$ then the argument is symmetric to above.
        \item If $e(\vec x)$ is $\mu x p(x,\vec x)$ then:
        \[
        \begin{array}{rll}
             & \mu x \id p(x,\vec f) + \mu x \id p(x,\vec f) 
             % & \text{by definition of $\id \cdot$}
             \\
              =  & \id p (\mu x \id p(x,\vec f),\vec f) + \id p (\mu x \id p(x,\vec f), \vec f) & \text{by $\mu$-unfolding, axiom \eqref{item:fixed-point-axioms}} \\
              = & \id p (\mu x \id p(x,\vec f),\vec f) & \text{by inductive hypothesis for $p(x,\vec x)$} \qedhere
        \end{array}
        \]
    \end{itemize}
\end{proof}

Note that, the above result notwithstanding, $\id e$ is typically not productive, even if $e$ is. In particular $\id a$ is not productive, even though $a$ is.

% \begin{remark}
% [The problem of interpreting $\idem$]
% \label{problem-of-interpreting-idem}
    Notice that, \cref{idem-trans-is-an-interp} notwithstanding, $\id \cdot$ does not necessarily interpret $\basicaxidfree + \idem$ into $\basicaxidfree$. This is because, a priori, a proof in $\basicaxidfree+\idem$ may exploit idempotency of expressions not in the image of $\id \cdot $ on productive expressions.
    % Perhaps proof theoretic methods could help prove an appropriate such result, but this is beyond the scope of the present work.
    For this reason, in what follows,
    we shall exploit the fact that we carefully kept track of the use of idempotency in our arguments in the previous section.
% \end{remark}

\subsection{Undecidability of $\basicaxidfree$ (and extensions)}
Recall $\top \df \mu x\left(\sum \Alphabet + \sum\limits_{a\in \Alphabet} ax\right)$.
By a similar argument to \cref{reg-tot-provable-with-idem} in the previous section we can obtain:

\begin{proposition}
\label{reg-tot-prov-under-idem-trans}
    Let $G \in \regtot$ and let $e_X$ be productive solutions of each of its non-terminals $X$ in $\basicaxidfree$, cf.~\cref{prod-cfg-has-prod-solns}, with $e_T = \top$. 
    Then $\basicaxidfree \proves \id\top \leq \id{e_S}$. 
\end{proposition}

For this we need appropriate versions of the two intermediate results we employed, \cref{lem:membership,lem:unfolding-top}, in the previous section.
Their proofs are morally similar to the previous ones, but we must be careful to state them appropriately in the absence of idempotency, exploiting the refined results of this section on productive CFGs and $\basicaxidfree$
\begin{lemma}
\label{lem:membership-productive}
    Fix a productive CFG $G$ and write $e_X$ for (any) corresponding productive solutions in $\basicaxidfree$ of each non-terminal $X$ of $G$, cf.~\cref{prod-cfg-has-prod-solns}.
    If $\vec a \in \lang X$ in $G$ then $\basicaxidfree \proves \id{\vec a} \leq \id e_X$.
\end{lemma}
% The proof of this is similar to that of \cref{lem:membership}, only relying on idempotency of $\id {\vec a}$ in place of $\vec a$, and using the refined results for productive CFGs and $\basicaxidfree$ in this section.
\begin{proof}
    By induction on the derivation tree of $\vec a $ from $X$ in $G$. Working in $\basicaxidfree$:
    \begin{itemize}
        \item If $X\to \vec a $ then we have $\vec a + g = e_X$ for some $g$ by \cref{prod-cfg-has-prod-solns}.
        Thus $\id{\vec a} + \id g = \id e_X$ by \cref{idem-trans-is-an-interp}, and so $\id{\vec a} \leq \id{e_X}$ by idempotency of $\id{\vec a}$, \cref{idem-trans-is-idem-on-prod-exprs}.
        \item The inductive step is the same as in the proof of \cref{lem:membership}. \qedhere
    \end{itemize}
\end{proof}


\begin{lemma}
\label{lem:unfolding-top-no-id}
    For $n>0$, we have $\basicaxidfree \proves \id\top = \sum\limits_{|\vec a|=1}^{n}\id{\vec a} + \sum\limits_{|\vec a|= n} \id{\vec a} \id\top$.
\end{lemma}
\begin{proof}
    The proof of \cref{lem:unfolding-top} does not use any identity axioms, and so goes through already in $\basicaxidfree$.
    The result now follows by \cref{idem-trans-is-an-interp}.
\end{proof}

We are now ready to prove the result we stated at the beginning of this section:

\begin{proof}
    [Proof of \cref{reg-tot-prov-under-idem-trans}]
    The proof goes through just like the proof of \cref{reg-tot-provable-with-idem}, only replacing all (productive) expressions by their image under $\id \cdot$.
    All appeals to \cref{lem:membership,lem:unfolding-top} are replaced by \cref{lem:membership-productive,lem:unfolding-top-no-id} respectively.
    Note that, by inspection of the proof, the only uses of idempotency (marked in \purple{purple}) are on productive expressions, whose image under $\id \cdot$ is indeed provably idempotent in $\basicaxidfree$ by \cref{idem-trans-is-idem-on-prod-exprs}.
    In particular, the expressions $e_Y$ and $e_S$ are productive by assumption.
\end{proof}

\subsection{Putting it together}
Finally we arrive at our promised strengthening of \cref{thm:undecidability-with-idem}:
\begin{theorem}
\label{thm:undecidability-without-idem}
    Any r.e.\ extension of $\basicaxidfree$ sound for $\Lang$ is $\Sigma^0_1$-complete.
\end{theorem}
The proof is similar to that of \cref{thm:undecidability-with-idem}, only being more careful about handling the $\id\cdot$ translation.
\begin{proof}
    Like in the proof of \cref{thm:undecidability-with-idem}, let us assume that all CFGs have start symbol $S$ and write $e^G_X$ for the productive solutions of nonterminals $X$ in a productive grammar $G$ obtained by \cref{prod-cfg-has-prod-solns}.
    Let $\mathsf U \supseteq \basicaxidfree$ be r.e.\ with $\Lang \models \mathsf U$.
    Now we write $\total{\mathsf U} \df \{G \text{ productive}: \mathsf U \proves \id \top \leq \id{(e^G_S)}\}$. We have:
    \begin{itemize}
        \item $\total {\mathsf U} \cap \nottot = \emptyset$, by assumption that $\Lang \models \mathsf U $ and by \cref{idem-trans-equiv-in-lang}.
        \item $\regtot \subseteq \total{\mathsf U}$, since $\basicaxidfree\subseteq \mathsf U$ and by \cref{reg-tot-prov-under-idem-trans}.
    \end{itemize}
    Thus by \cref{separator-of-eff-insep-sets-is-S01-complete,regtot-nottot-eff-insep} we have that $\total{\mathsf U}$ is $\Sigma^0_1$-complete.
    Again, by the obvious computable reduction from $\total{\mathsf U}$ to $\mathsf U$, we conclude as required. 
\end{proof}


\begin{corollary}
    The equational theories of Conway $\mu$-semirings and Park $\mu$-semirings are undecidable, and moreover $\Sigma^0_1$-complete.
    The same holds even for identity-free versions of these theories.
\end{corollary}

% \section{An interpretation of idempotency in the presence of induction}

% Recalling \cref{problem-of-interpreting-idem}, we cannot easily interpret idempotency in non-idempotent theories, despite the technology built up in the previous section, requiring a more low-level analysis of idempotency.
% In this section we point out that a bona fide interpretation becomes available in the presence of induction:

% \begin{theorem}
%     $\basicaxidfree + \idem + \ind$ is interpreted by $\basicaxidfree + \ind$ under $\id \cdot$.\todo{statement must make clear that the target theory is equipped with a partial order, not necessarily the natural one.}
% \end{theorem}
% Note, in this statement 
% Recalling the proof of \cref{idem-trans-is-an-interp}, for the above result it suffices to show that the image of $\id \cdot $ is idempotent on \emph{all} expressions, not just the productive ones, cf.~\cref{idem-trans-is-idem-on-prod-exprs}.
%     Naturally for this we must exploit induction, resolving `unguarded' bound variables to $0$.


% \todo[inline]{write this proof}

% \begin{corollary}
%     The theory of naturally ordered idempotent identity-free Park $\mu$-semirings is interpreted by the the theory of (not necessarily idempotent) identity-free Park $\mu$-semirings.
% \end{corollary}

% \subsection{New}

% Recall $e' := \mu y (e + \mu z (y + z))$ is provably idempotent in $\basicaxidfree$.
% \begin{definition}
%     [$j$-interpretation]
%     For \emph{all} $1$-free expressions 
% \end{definition}


% \section{Further consequences}



% Working in a setting without idempotency, let us point out that an expression $e$ has an idempotent upper bound $e'\df \mu X (e + \mu Y(X+Y))$.
% In fact, under idempotency, we have $e = e'$.
% On the other hand under induction $e'$ is the \emph{least} idempotent upper bound of $e$.

% \todo[inline]{
% If we replace all letters $a$ in a CFG by $a'$, and the CFG only has producing transitions, this should force idempotency of solutions. This is beacuse, recall the class of idempotents forms an ideal. Still need to deal with $1$ though.
% }
